\scanpage{16}
\mspage{\unsure{160|16a}}

algébriquement clos. Mais alors, $s$ désignant le point fermé de $S$, comme $G$ est étale et séparé sur $S$, il contient un sous-schéma à la fois ouvert et fermé $G'$, \inserted{fini sur $S$,} tel que $G'_s=G_s$ \inserted{(EGA II 6.2.6.)}, et on voit de suite que $G'$ est ici un sous-groupe de $G$. \deleted{\unsure{Par suite la fibre géométrique $f(s)$}} D'ailleurs $G'$ \deleted{qui} étant étale fini sur $S$ \deleted{qui} $=\operatorname{Spec}(V)$, $V$ complet à corps résiduel algébriquement clos, est un groupe constant, donc de la forme $A_S$, où $A=G'(k(s))=G(k(s))$ a pour classe $f(s)$. \deleted{On voit donc} Si $B$ est la fibre géométrique de $G$ en le point générique $t$ de $S$, on a donc un \deleted{\unsure{homomor}}\inserted{monomor}phisme canonique $A\longrightarrow B$, ce qui prouve $f(s)\leq f(t)$. (NB cette démonstration prouve en fait la semi-continuité pour une relation d'ordre sur $E$ plus fine que celle indiquée dans 2.1.). Le fait que $f$ soit continue en $s$ si et seulement si $s\rightsquigarrow\operatorname{card}f(s)$ l'est résulte du fait que pour $w,w'\in E$, les relations \deleted{\unsure{w}} $w\leq w'$ et $\operatorname{card}w=\operatorname{card}w'$ impliquent $w=w'$. Le fait que cette condition soit équivalente à \deleted{\unsure{à}} la finitude de $W$ sur un voisinage de $s$ est alors indépendant de\deleted{s}\inserted{la} structure\deleted{s} de groupe\deleted{s}\inserted{sur $W$}, et a été signalé après \inserted{Exp XI} 5.10. ; d'ailleurs sa démonstration \deleted{qui} se fait aisément par les arguments précédents, \inserted{en} utilisant le critère valuatif de propreté EGA II 7.3.8.

\uline{Lemme 2.3.} Soient $G$ un groupe algébrique affine lisse sur un corps algébriquement clos $k$, $R\subset T$ deux sous-tores, $W(R)$ et $W(T)$ les deux groupes finis associés comme dans \inserted{Exp XI} 5.9., quotients du normalisateur par le centralisateur. Alors $W(R)$ est isomorphe à un quotient d'un sous-groupe de $W(T)$.

\deleted{Dans la démonstration qui suit, il sera commode d'interpréter $G$, $R,T$, $C(R),N(R)$, $C(T)$, $N(T)$ de façon purement ensembliste, comme les groupes de points des préschémas $G,R,T$ dans $k$, et les centralisateurs et nor}

\mspage{15}
